\section{Local heights and the two external inputs}\label{sec:inputs}

We isolate the source-dependent part of the argument. All subsequent
orbit, field and parameter-infinitude calculations are proved below.

\begin{proposition}[Lee--Nam, specialized reduction]\label{prop:reduction}
If $|\Prep(f;a,b)|=\infty$ and $a,b$ are not both in $k$, then
\eqref{eq:three} holds.
\end{proposition}

This is the specialization of \citet[Theorem~1.5, proved as
Theorem~4.4]{LeeNam2025}. Both prime-to-three degree parts exceed one,
and the weights are equal as calculated above. In the nonzero case
their relation reads
\[
 (f(b)-f(a))^{\,3^1-3^0}=-\frac{4}{2}=1
 \quad\text{in characteristic three},
\]
whose roots are $1$ and $-1$. We use the cited theorem for the
necessity of this relation, not just the elementary solution of it.

For the second input, let $K$ be a finite extension of the perfect
closure of $k(t)$. A place $v$ extends a function-field place trivial
on $k$; write $|\cdot|_v$ for a compatible nonarchimedean absolute
value. It extends to an algebraic closure and its completion
$\mathbb C_v$. For $x,\lambda\in\mathbb C_v$ the local canonical
height for our degree-six polynomial is
\begin{equation}\label{eq:height}
 \hhat_{v,\lambda}(x)
 =\lim_{n\to\infty}6^{-n}\log^+|F_\lambda^n(x)|_v,\qquad
 \log^+u=\log\max\{1,u\}.
\end{equation}
In particular any point with finite forward orbit has local height zero.

\begin{proposition}[Universal local-height equality]\label{prop:heights}
Let $a,b\in K$, with $K$ as above. If infinitely many distinct
$\lambda\in\overline K$ make both $a$ and $b$ preperiodic under
$F_\lambda$, then
\begin{equation}\label{eq:equalheight}
 \hhat_{v,\lambda}(a)=\hhat_{v,\lambda}(b)
 \quad\text{for every place }v
 \text{ and every }\lambda\in\mathbb C_v .
\end{equation}
\end{proposition}

We use \citet[Theorem~2.1]{LeeNam2025}, stated there from
Ghioca--Hsia \citep[Theorem~2.13]{GhiocaHsia2026}.
The polynomial conditions for that input are monicity, degree at
least two and $f(0)=0$, all satisfied here. There is no strict
exponent-weight hypothesis in this height statement. It asserts
equality at \emph{every} parameter, not only at the original
common-preperiodicity parameters. That quantifier is essential:
our separating parameter will contradict the equality, not belong
to the common-parameter set.

No discreteness of the extended value group will be needed. At a
place over the pole of a transcendental $a$, nonzero elements of $k$
have absolute value one and $|a|_v>1$. These facts suffice for the
strict dominance estimates below.
